% ============================================================
% 第二章 数学基础与差分格式
% ============================================================
\section{数学基础与差分格式}
\label{sec:math}

\subsection{Poisson 方程与 Helmholtz 方程}

考虑矩形区域 $\Omega = [0, L_x] \times [0, L_y]$ 上的椭圆型方程：
\begin{equation}
  \mathcal{L}u \equiv (-\nabla^2 + \sigma)u = f, \quad (x,y) \in \Omega
  \label{eq:helmholtz}
\end{equation}
其中 $\sigma$ 为参数，其取值决定方程类型：
\begin{itemize}[itemsep=2pt]
  \item $\sigma = 0$：Poisson 方程 $-\nabla^2 u = f$；
  \item $\sigma = +\kappa^2$（$\kappa > 0$）：modified Helmholtz 方程（screened Poisson 方程），
        离散矩阵对称正定，无共振；
  \item $\sigma = -\kappa^2$：true Helmholtz 方程，离散矩阵可能不定，并可能发生共振。
\end{itemize}
本文在统一框架~\eqref{eq:helmholtz}下研究三类方程的快速数值求解。

常见的边界条件包括：
\begin{itemize}[itemsep=2pt]
  \item \textbf{Dirichlet 边界条件}：$u|_{\partial\Omega} = g_D$
  \item \textbf{Neumann 边界条件}：$\frac{\partial u}{\partial n}\big|_{\partial\Omega} = g_N$
\end{itemize}

\subsection{五点差分格式（二阶精度）}

对区域 $\Omega$ 进行均匀网格剖分，取步长 $h = L_x / (n-1)$，
节点 $(x_i, y_j) = (ih, jh)$，$i,j = 0, 1, \ldots, n-1$。

\subsubsection{Taylor 展开推导}

对 $u(x_{i+1}, y_j)$ 和 $u(x_{i-1}, y_j)$ 分别进行 Taylor 展开：
\begin{align}
  u(x_{i+1}, y_j) &= u_{i,j} + h\frac{\partial u}{\partial x}\bigg|_{i,j}
  + \frac{h^2}{2}\frac{\partial^2 u}{\partial x^2}\bigg|_{i,j}
  + \frac{h^3}{6}\frac{\partial^3 u}{\partial x^3}\bigg|_{i,j}
  + \frac{h^4}{24}\frac{\partial^4 u}{\partial x^4}\bigg|_{i,j} + O(h^5) \\
  u(x_{i-1}, y_j) &= u_{i,j} - h\frac{\partial u}{\partial x}\bigg|_{i,j}
  + \frac{h^2}{2}\frac{\partial^2 u}{\partial x^2}\bigg|_{i,j}
  - \frac{h^3}{6}\frac{\partial^3 u}{\partial x^3}\bigg|_{i,j}
  + \frac{h^4}{24}\frac{\partial^4 u}{\partial x^4}\bigg|_{i,j} + O(h^5)
\end{align}

两式相加消去奇数阶项，得到二阶导数的中心差分展开：
\begin{equation}
  \frac{u_{i-1,j} - 2u_{i,j} + u_{i+1,j}}{h^2}
  =
  \frac{\partial^2 u}{\partial x^2}\bigg|_{i,j}
  + \frac{h^2}{12}\frac{\partial^4 u}{\partial x^4}\bigg|_{i,j} + O(h^4)
  \label{eq:central_diff_x}
\end{equation}

同理对 $y$ 方向有：
\begin{equation}
  \frac{u_{i,j-1} - 2u_{i,j} + u_{i,j+1}}{h^2}
  =
  \frac{\partial^2 u}{\partial y^2}\bigg|_{i,j}
  + \frac{h^2}{12}\frac{\partial^4 u}{\partial y^4}\bigg|_{i,j} + O(h^4)
  \label{eq:central_diff_y}
\end{equation}

将~\eqref{eq:central_diff_x}和~\eqref{eq:central_diff_y}代入 $-\nabla^2 u = -\frac{\partial^2 u}{\partial x^2} - \frac{\partial^2 u}{\partial y^2}$，得到五点差分格式：
\begin{equation}
  \frac{1}{h^2}(4u_{i,j} - u_{i-1,j} - u_{i+1,j} - u_{i,j-1} - u_{i,j+1}) + \sigma u_{i,j} = f_{i,j}
  \label{eq:5pt}
\end{equation}

记五点正 Laplacian 离散算子为
\[
L_h^{(5)}u_{i,j}
=
\frac{4u_{i,j}-u_{i-1,j}-u_{i+1,j}-u_{i,j-1}-u_{i,j+1}}{h^2}.
\]
该算子近似的是 $-\Delta u$，而不是 $\Delta u$。由上述 Taylor 展开可得
\[
L_h^{(5)}u_{i,j}
=
-\Delta u(x_i,y_j)
-
\frac{h^2}{12}
\left(
\frac{\partial^4 u}{\partial x^4}
+
\frac{\partial^4 u}{\partial y^4}
\right)_{i,j}
+O(h^4).
\]

本文在本节采用 residual 方向
$\tau_h(u)=\bigl(L_h^{(5)}+\sigma I\bigr)u-f$，
即将精确解代入离散算子后减去连续方程右端。
在该约定下，五点格式的截断误差为：
\begin{equation}
  \tau_{i,j} = -\frac{h^2}{12}\left(\frac{\partial^4 u}{\partial x^4} + \frac{\partial^4 u}{\partial y^4}\right)\bigg|_{i,j} + O(h^4) = O(h^2)
  \label{eq:truncation_5pt}
\end{equation}

\begin{theorem}[五点差分格式的截断误差]
  \label{thm:5pt_truncation}
  对于充分光滑的解 $u$，五点差分格式~\eqref{eq:5pt} 的截断误差为 $O(h^2)$。
  当 $u$ 为调和函数（$\nabla^2 u = 0$）时，若 $\frac{\partial^4 u}{\partial x^4} + \frac{\partial^4 u}{\partial y^4} = 0$，
  则截断误差可达 $O(h^4)$（超收敛现象）。
\end{theorem}

\subsubsection{矩阵表示}

将内节点按行优先排列为向量 $\mathbf{u} \in \mathbb{R}^{(n-2)^2}$，
可得线性方程组：
\begin{equation}
  A\mathbf{u} = \mathbf{f}
  \label{eq:linear-system}
\end{equation}
其中系数矩阵 $A$ 利用 Kronecker 积表示为：
\begin{equation}
  A = \frac{1}{h^2}(I_{N} \otimes B + B \otimes I_{N}) + \sigma I_{N^2}
  \label{eq:matrix_A_kron}
\end{equation}
$B$ 为 $(n-2) \times (n-2)$ 一维 Laplacian 矩阵：
\begin{equation}
  B = \begin{pmatrix}
    2 & -1 & & \\
    -1 & 2 & \ddots & \\
    & \ddots & \ddots & -1 \\
    & & -1 & 2
  \end{pmatrix}
  \label{eq:B_matrix}
\end{equation}
$N = n-2$ 为内节点数。
采用分解~\eqref{eq:matrix_A_kron}的优势在于：
一维特征值 $\mu_p = 2-2\cos(p\pi/(N+1))$ 简洁无歧义，
二维特征值直接叠加为 $(\mu_p+\mu_q)/h^2+\sigma$，避免了 double counting 问题。

\subsubsection{特征值与特征向量的推导}

\begin{theorem}[DST-I 对角化与特征值]
  \label{thm:eigenvalues}
  矩阵 $B$ 可被 DST-I 矩阵对角化。令 $N = n-2$，定义 DST-I 变换矩阵
  $S \in \mathbb{R}^{N \times N}$，其元素为：
  \begin{equation}
    S_{p,i} = \sqrt{\frac{2}{N+1}} \sin\frac{p i \pi}{N+1}, \quad p,i = 1, \ldots, N
  \end{equation}
  则 $S$ 为正交矩阵（$S^T S = I$），且
  \begin{equation}
    B = S^T \Lambda_B S, \quad \Lambda_B = \mathrm{diag}(\mu_1, \ldots, \mu_N)
  \end{equation}
  一维特征值为：
  \begin{equation}
    \mu_p = 2 - 2\cos\frac{p\pi}{N+1}, \quad p = 1, 2, \ldots, N
    \label{eq:mu_p}
  \end{equation}

  \textbf{证明}：验证 $\mathbf{v}^{(p)} = (\sin\frac{p\pi}{N+1}, \sin\frac{2p\pi}{N+1}, \ldots, \sin\frac{Np\pi}{N+1})^T$
  是 $B$ 的特征向量。计算 $B\mathbf{v}^{(p)}$ 的第 $i$ 个分量：
  \begin{equation}
    (B\mathbf{v}^{(p)})_i = 2\sin\frac{ip\pi}{N+1} - \sin\frac{(i-1)p\pi}{N+1} - \sin\frac{(i+1)p\pi}{N+1}
  \end{equation}
  利用和差化积公式 $\sin\alpha + \sin\beta = 2\sin\frac{\alpha+\beta}{2}\cos\frac{\alpha-\beta}{2}$：
  \begin{equation}
    \sin\frac{(i-1)p\pi}{N+1} + \sin\frac{(i+1)p\pi}{N+1}
    = 2\sin\frac{ip\pi}{N+1}\cos\frac{p\pi}{N+1}
  \end{equation}
  因此：
  \begin{equation}
    (B\mathbf{v}^{(p)})_i = (2 - 2\cos\frac{p\pi}{N+1})\sin\frac{ip\pi}{N+1}
    = \mu_p \cdot v_i^{(p)} \qquad \qed
  \end{equation}
\end{theorem}

对于二维问题，由 Kronecker 积~\eqref{eq:matrix_A_kron}的性质，矩阵 $A$ 的特征值为：
\begin{equation}
  d_{p,q} = \frac{\mu_p + \mu_q}{h^2} + \sigma
  = \frac{2}{h^2}\left(2 - \cos\frac{p\pi}{N+1} - \cos\frac{q\pi}{N+1}\right) + \sigma
  \label{eq:2d_eigenvalue}
\end{equation}
对应的特征向量为 $\mathbf{v}^{(p)} \otimes \mathbf{v}^{(q)}$（Kronecker 积），
即 $v_{i,j}^{(p,q)} = \sin\frac{ip\pi}{N+1}\sin\frac{jq\pi}{N+1}$。
频域分母统一为 $d_{p,q} = \lambda_{p,q} + \sigma$，
其中 $\lambda_{p,q} = (\mu_p+\mu_q)/h^2$ 为 Poisson 方程（$\sigma=0$）的特征值。

\begin{corollary}[条件数估计]
  \label{cor:condition_number}
  五点差分矩阵 $A$（$\sigma = 0$，Dirichlet BC）的条件数为：
  \begin{equation}
    \kappa(A) = \frac{d_{\max}}{d_{\min}}
    \sim \frac{8/h^2}{2\pi^2}
    = \frac{4}{\pi^2 h^2} = O(h^{-2})
    \label{eq:condition_number}
  \end{equation}
  其中 $\lambda_{\min} = d_{1,1} \sim 2\pi^2$（leading-order 渐近，精确值为 $8\sin^2(\pi h/2)/h^2$），
  $\lambda_{\max} = d_{N,N} \sim 8/h^2$。
  条件数随网格加密以 $O(h^{-2})$ 增长，这是椭圆型方程差分离散的一般规律。

  对于固定网格和同一 Dirichlet 五点 SPD 离散系统，当 $\kappa>0$ 时，
  modified Helmholtz 方程（$\sigma = +\kappa^2$）的 2-范数条件数为：
  \begin{equation}
    \kappa_2(A_{\mathrm{mod}})
    =
    \frac{\lambda_{\max}+\kappa^2}{\lambda_{\min}+\kappa^2}
    <
    \frac{\lambda_{\max}}{\lambda_{\min}}
    =
    \kappa_2(A_{\mathrm{Poisson}})
    \label{eq:condition_modified}
  \end{equation}
  该严格不等式依赖 $0<\lambda_{\min}<\lambda_{\max}$ 与 $\kappa>0$；
  若 $\kappa=0$，modified Helmholtz 退化为 Poisson，二者条件数相等。

  对于 true Helmholtz 方程（$\sigma = -\kappa^2$），
  当 $\kappa^2$ 接近某特征值 $\lambda_{p,q}$ 时，分母 $d_{p,q}=\lambda_{p,q}-\kappa^2$ 接近零，
  谱分母指标趋于无穷，即 near-resonance 小分母现象。
  此时矩阵可能接近奇异、不定，或在 $\kappa^2>\lambda_{\max}$ 时负定，不能按 SPD 条件数解释。
\end{corollary}

\subsection{九点紧致差分格式（FFT9）}

标准五点差分格式的截断误差主项为 $O(h^2)$。
通过引入紧致差分思想\cite{Lynch1977}——在格式中同时使用 $R_h$ 算子（质量矩阵算子）——
可以在不扩大计算模板的前提下将精度提升至 $O(h^4)$。

\subsubsection{$R_h$ 算子的定义与 Taylor 展开}

定义算子 $R_h$（质量矩阵算子）：
\begin{equation}
  R_h v_{i,j} = \frac{1}{12}(v_{i-1,j} + v_{i+1,j} + v_{i,j-1} + v_{i,j+1})
  + \frac{2}{3}v_{i,j}
  \label{eq:Rh_def}
\end{equation}

$R_h$ 的物理意义是对场变量进行加权平均，中心权重为 $2/3$，
四个邻点权重各为 $1/12$。以下推导 $R_h$ 的 Taylor 展开式。

\begin{remark}[$R_h$ 在边界附近的取值]
  \label{rem:Rh_boundary}
  对于 Dirichlet 边界条件，$R_h$ 的作用域仅限于内节点 $(i,j)$，$1 \leq i,j \leq n-2$。
  当内节点邻接边界时（如 $i=1$），$R_h$ 模板中的 $v_{0,j}$ 为已知边界值，
  需移至右端项处理（详见第~\ref{sec:fft}章的边界修正讨论）。
  对于 Neumann 边界条件，$R_h$ 的模板涉及边界节点，
  其中邻接边界的项需利用 ghost-point 关系 $v_{-1,j} = v_{1,j}$ 替换，
  这导致 $R_h$ 在边界行/列的权重与内点不同。
  在频域求解中，上述边界效应通过物理空间的边界修正项统一处理，
  不改变 FFT9 原始频域分母 $D_{p,q}^{\mathrm{raw}}$ 的结构。
\end{remark}

\begin{lemma}[$R_h$ 的 Taylor 展开]
  \label{lem:Rh_taylor}
  对充分光滑的函数 $v$，算子 $R_h$ 满足：
  \begin{equation}
    R_h v = v + \frac{h^2}{12}\nabla^2 v + \frac{h^4}{144}
    \left(\frac{\partial^4 v}{\partial x^4} + \frac{\partial^4 v}{\partial y^4}\right)
    + O(h^6)
    \label{eq:Rh_taylor}
  \end{equation}

  \textbf{证明}：对 $v_{i\pm1,j}$ 和 $v_{i,j\pm1}$ 做 Taylor 展开：
  \begin{align}
    v_{i\pm1,j} &= v \pm h\frac{\partial v}{\partial x} + \frac{h^2}{2}\frac{\partial^2 v}{\partial x^2}
    \pm \frac{h^3}{6}\frac{\partial^3 v}{\partial x^3} + \frac{h^4}{24}\frac{\partial^4 v}{\partial x^4} + O(h^5) \\
    v_{i,j\pm1} &= v \pm h\frac{\partial v}{\partial y} + \frac{h^2}{2}\frac{\partial^2 v}{\partial y^2}
    \pm \frac{h^3}{6}\frac{\partial^3 v}{\partial y^3} + \frac{h^4}{24}\frac{\partial^4 v}{\partial y^4} + O(h^5)
  \end{align}
  其中所有偏导数在 $(x_i, y_j)$ 处取值。将四式求和，奇数阶项消去：
  \begin{equation}
    v_{i-1,j} + v_{i+1,j} + v_{i,j-1} + v_{i,j+1}
    = 4v + h^2\nabla^2 v + \frac{h^4}{12}\left(\frac{\partial^4 v}{\partial x^4}
    + \frac{\partial^4 v}{\partial y^4}\right) + O(h^6)
  \end{equation}
  代入~\eqref{eq:Rh_def}：
  \begin{equation}
    \begin{aligned}
    R_h v
    &= \frac{1}{12}\left[4v + h^2\nabla^2 v
       + \frac{h^4}{12}\left(\frac{\partial^4 v}{\partial x^4}
       + \frac{\partial^4 v}{\partial y^4}\right)\right]
       + \frac{2}{3}v + O(h^6) \\
    &= v + \frac{h^2}{12}\nabla^2 v
       + \frac{h^4}{144}\left(\frac{\partial^4 v}{\partial x^4}
       + \frac{\partial^4 v}{\partial y^4}\right) + O(h^6). \qquad \qed
    \end{aligned}
  \end{equation}
\end{lemma}

\subsubsection{紧致差分格式的推导}

利用引理~\ref{lem:Rh_taylor}，将 $R_h$ 作用于方程 $(-\nabla^2 + \sigma)u = f$ 的两端：
\begin{equation}
  R_h(-\nabla^2 u) + \sigma R_h u = R_h f
  \label{eq:Rh_apply}
\end{equation}

由~\eqref{eq:Rh_taylor}，$R_h u = u + O(h^2)$，因此 $\sigma R_h u = \sigma u + O(h^2)$。
关键在于 $R_h(-\nabla^2 u)$ 的展开。令 $w = -\nabla^2 u$，利用~\eqref{eq:Rh_taylor}：
\begin{equation}
  R_h w = w + \frac{h^2}{12}\nabla^2 w + O(h^4)
  = -\nabla^2 u + \frac{h^2}{12}\nabla^2(-\nabla^2 u) + O(h^4)
  = -\nabla^2 u - \frac{h^2}{12}\nabla^4 u + O(h^4)
  \label{eq:Rh_nabla2}
\end{equation}

另一方面，五点正 Laplacian 算子 $L_h^{(5)}$（即 $-\Delta$ 的离散化）的截断误差为：
\begin{equation}
  L_h^{(5)} u = -\nabla^2 u
  - \frac{h^2}{12}\left(\frac{\partial^4 u}{\partial x^4} + \frac{\partial^4 u}{\partial y^4}\right) + O(h^4)
  \label{eq:L5_error}
\end{equation}

对比~\eqref{eq:Rh_nabla2}和~\eqref{eq:L5_error}：
$R_h(-\nabla^2 u)$ 中的 $O(h^2)$ 误差为 $-\frac{h^2}{12}\nabla^4 u$，
而 $L_h^{(5)} u$ 中的 $O(h^2)$ 误差为 $-\frac{h^2}{12}(\frac{\partial^4 u}{\partial x^4} + \frac{\partial^4 u}{\partial y^4})$。
注意到 $\nabla^4 u = \frac{\partial^4 u}{\partial x^4} + 2\frac{\partial^4 u}{\partial x^2 \partial y^2} + \frac{\partial^4 u}{\partial y^4}$，
两者的差恰好为混合偏导数项：
\begin{equation}
  R_h(-\nabla^2 u) - L_h^{(5)} u = -\frac{h^2}{12}\left(\nabla^4 u
  - \frac{\partial^4 u}{\partial x^4} - \frac{\partial^4 u}{\partial y^4}\right) + O(h^4)
  = -\frac{h^2}{6}\frac{\partial^4 u}{\partial x^2 \partial y^2} + O(h^4)
  \label{eq:diff_Rh_L5}
\end{equation}

这表明：若能在 $L_h^{(5)}$ 的基础上减去 $\frac{h^2}{6}\frac{\partial^4 u}{\partial x^2 \partial y^2}$ 项，
则可消去 $O(h^2)$ 误差，达到 $O(h^4)$ 精度。

混合偏导数 $\frac{\partial^4 u}{\partial x^2 \partial y^2}$ 的标准差分近似为：
\begin{equation}
  \begin{aligned}
  \frac{\partial^4 u}{\partial x^2 \partial y^2}\bigg|_{i,j}
  = \frac{1}{h^4}\big(&u_{i+1,j+1} + u_{i-1,j+1}
  + u_{i+1,j-1} + u_{i-1,j-1} \\
  &- 2u_{i+1,j} - 2u_{i-1,j}
  - 2u_{i,j+1} - 2u_{i,j-1} + 4u_{i,j}\big)
  + O(h^2)
  \end{aligned}
  \label{eq:mixed_deriv}
\end{equation}

将~\eqref{eq:mixed_deriv}代入~\eqref{eq:diff_Rh_L5}，与 $L_h^{(5)}$ 合并，
经化简得到九点 Laplacian 算子 $L_h$，其相反数近似 $R_h(-\nabla^2)$：
\begin{equation}
  \begin{aligned}
  L_h v_{i,j}
  = \frac{1}{6h^2}\big[&
  4(v_{i-1,j} + v_{i+1,j} + v_{i,j-1} + v_{i,j+1}) \\
  &+ (v_{i-1,j-1} + v_{i-1,j+1}
      + v_{i+1,j-1} + v_{i+1,j+1})
  - 20v_{i,j}\big]
  \end{aligned}
  \label{eq:Lh_9pt}
\end{equation}
可以验证 $-L_h u = R_h(-\nabla^2 u) + O(h^4)$，即九点模板的相反数与 $R_h$ 作用后的精确算子
仅在 $O(h^4)$ 级别存在偏差。

\begin{lemma}[$L_h$ 的算子精确分解]
  \label{lem:Lh_decompose}
  若 $L_h^{(5)}$ 表示五点正 Laplacian（即 $-\Delta$ 的离散化），则九点 Laplacian 算子满足：
  \begin{equation}
    -L_h = L_h^{(5)} - \frac{h^2}{6}C_h
    \label{eq:Lh_decompose}
  \end{equation}
  其中交叉模板 $C_h$ 定义为：
  \begin{equation}
    C_h v_{i,j} = \frac{1}{h^4}(v_{i+1,j+1} + v_{i-1,j+1} + v_{i+1,j-1} + v_{i-1,j-1}
    - 2v_{i+1,j} - 2v_{i-1,j} - 2v_{i,j+1} - 2v_{i,j-1} + 4v_{i,j})
    \label{eq:Ch_def}
  \end{equation}
  $C_h$ 恰好是混合偏导数 $\partial^4/\partial x^2\partial y^2$ 的标准差分近似，
  即 $C_h v = \frac{\partial^4 v}{\partial x^2\partial y^2} + O(h^2)$。

  \textbf{证明}：将~\eqref{eq:Lh_9pt}中的 $-L_h$ 模板与五点正 Laplacian 模板~\eqref{eq:5pt}做差，可得
  \begin{align}
    L_h^{(5)}v - (-L_h v)
    &= \frac{1}{6h^2}\Big[
    v_{i+1,j+1}+v_{i-1,j+1}+v_{i+1,j-1}+v_{i-1,j-1} \notag\\
    &\quad -2v_{i+1,j}-2v_{i-1,j}-2v_{i,j+1}-2v_{i,j-1}+4v_{i,j}\Big]
    = \frac{h^2}{6}C_h v .
  \end{align}
  因此 $-L_h = L_h^{(5)}-\frac{h^2}{6}C_h$。
  由~\eqref{eq:mixed_deriv}，$C_h$ 是 $\partial^4/\partial x^2\partial y^2$ 的差分近似。$\qed$
\end{lemma}

\begin{remark}
  分解~\eqref{eq:Lh_decompose}的物理意义十分清晰：$L_h$ 相比 $L_h^{(5)}$ 多出的部分
  恰好是对混合偏导数 $\partial^4 u/\partial x^2\partial y^2$ 的补偿。
  这正是~\eqref{eq:diff_Rh_L5}中需要消去的项。
  因此紧致差分格式可以理解为：先用 $R_h$ 识别出 $O(h^2)$ 误差
  （混合偏导数项），再用四个角节点的差分模板精确补偿该项。
\end{remark}

最终得到紧致差分格式：
\begin{equation}
  (-L_h + \sigma R_h)u_{i,j} = R_h f_{i,j}
  \label{eq:compact-9pt}
\end{equation}
其中 $L_h \approx \Delta = \nabla^2$（非 $-\Delta$），因此 $-L_h \approx -\nabla^2$。
频域求解时，分母为 $-\hat{\lambda}_L + \sigma\hat{\lambda}_R$，
对 modified Helmholtz（$\sigma=+\kappa^2$）恒为正（因 $-\hat{\lambda}_L > 0$ 且 $\hat{\lambda}_R > 0$），
对 true Helmholtz（$\sigma=-\kappa^2$）可能接近零。

\begin{remark}
  格式~\eqref{eq:compact-9pt}中 $R_h$ 出现在等式两端，使得该格式本质上是
  一个\textbf{广义特征值问题}：$L_h$ 和 $R_h$ 同时作用于 $u$，
  右端项也需要用 $R_h$ 修正。这种"紧致"结构保证了使用仅九个节点的模板
  即可在内点 Fourier symbol 层面达到四阶一致性；
  当前 Dirichlet 实现的整体四阶收敛行为将在第~\ref{sec:experiments}章通过 manufactured solution 数值实验验证。
\end{remark}

\subsubsection{Fourier 特征值的推导}

为建立四阶一致性，采用 Fourier 分析方法。
设离散模式为 $u_{i,j} = e^{\mathrm{i}(\alpha i + \beta j)}$，
其中 $\alpha = p\pi h$，$\beta = q\pi h$（$h$ 为步长，$p,q$ 为模式指标）。

\begin{lemma}[Fourier 特征值]
  \label{lem:fourier_eigenvalues}
  算子 $L_h$ 和 $R_h$ 在 Fourier 模式下的特征值分别为：
  \begin{align}
    \hat{\lambda}_L(\alpha,\beta) &= \frac{1}{6h^2}(-20 + 8\cos\alpha + 8\cos\beta + 4\cos\alpha\cos\beta)
    \label{eq:lam_L} \\
    \hat{\lambda}_R(\alpha,\beta) &= \frac{2}{3} + \frac{1}{6}(\cos\alpha + \cos\beta)
    \label{eq:lam_R}
  \end{align}

  由于本文的九点算子 $L_h$ 近似的是连续 Laplacian $\Delta$，而不是 $-\Delta$，
  因此对 Dirichlet 非零模态有 $\hat{\lambda}_L(p,q)<0$。
  真正对应正算子 $-\Delta$ 的离散特征量是 $-\hat{\lambda}_L(p,q)>0$。
  同时，$R_h$ 的 Fourier symbol 满足 $\hat{\lambda}_R(p,q)>0$。

  \textbf{证明}：将 $u_{i,j} = e^{\mathrm{i}(\alpha i + \beta j)}$ 代入 $L_h$：
  \begin{align}
    L_h u_{i,j} &= \frac{1}{6h^2}\left[4(e^{-\mathrm{i}\alpha}+e^{\mathrm{i}\alpha}+e^{-\mathrm{i}\beta}+e^{\mathrm{i}\beta})
    + (e^{-\mathrm{i}\alpha-\mathrm{i}\beta}+e^{-\mathrm{i}\alpha+\mathrm{i}\beta}+e^{\mathrm{i}\alpha-\mathrm{i}\beta}+e^{\mathrm{i}\alpha+\mathrm{i}\beta})
    - 20\right] u_{i,j} \notag \\
    &= \frac{1}{6h^2}\left[8\cos\alpha + 8\cos\beta + 4\cos\alpha\cos\beta - 20\right] u_{i,j}
  \end{align}
  其中利用了 $e^{-\mathrm{i}\alpha}+e^{\mathrm{i}\alpha} = 2\cos\alpha$ 以及
  $(e^{-\mathrm{i}\alpha}+e^{\mathrm{i}\alpha})(e^{-\mathrm{i}\beta}+e^{\mathrm{i}\beta}) = 4\cos\alpha\cos\beta$。

  同理代入 $R_h$：
  \begin{equation}
    R_h u_{i,j} = \left[\frac{1}{12}(2\cos\alpha + 2\cos\beta) + \frac{2}{3}\right] u_{i,j}
    = \left[\frac{2}{3} + \frac{1}{6}(\cos\alpha + \cos\beta)\right] u_{i,j} \qquad \qed
  \end{equation}
\end{lemma}

\subsubsection{FFT9 Fourier symbol 四阶一致性分析}

\begin{theorem}[FFT9 紧致格式的 Fourier symbol 四阶一致性]
  \label{thm:4th_order}
  设
  \[
    \alpha=\xi h,\qquad \beta=\eta h,
  \]
  其中 $\xi,\eta$ 为固定波数。对于九点紧致格式
  \[
    (-L_h+\sigma R_h)u=R_hf,
  \]
  其有效离散特征值为
  \begin{equation}
    \lambda_{\mathrm{discrete}}(\alpha,\beta)
    =
    -\frac{\hat{\lambda}_L(\alpha,\beta)}
    {\hat{\lambda}_R(\alpha,\beta)}
    +\sigma.
  \end{equation}
  当 $h\to0$ 时，
  \begin{equation}
    \lambda_{\mathrm{discrete}}(\alpha,\beta)
    =
    \xi^2+\eta^2+\sigma
    -
    \frac{(\xi^2+\eta^2)(3\xi^4-8\xi^2\eta^2+3\eta^4)}{720}h^4
    +O(h^6).
    \label{eq:fft9_symbol_4th}
  \end{equation}
  因此 $h^2$ 项完全消去，FFT9 格式在内点 Fourier symbol 意义下对
  $(-\Delta+\sigma)$ 具有四阶一致性。

  \textbf{证明}：
  格式~\eqref{eq:compact-9pt} 在 Fourier 空间中为：
  \begin{equation}
    (-\hat{\lambda}_L + \sigma\hat{\lambda}_R)\hat{u} = \hat{\lambda}_R\hat{f}
    \quad \Longrightarrow \quad
    \hat{u} = \frac{\hat{\lambda}_R}{-\hat{\lambda}_L + \sigma\hat{\lambda}_R}\hat{f}
    = \frac{\hat{f}}{-\hat{\lambda}_L/\hat{\lambda}_R + \sigma}
  \end{equation}
  有效离散特征值为
  \[
    \lambda_{\mathrm{discrete}}
    =
    -\frac{\hat{\lambda}_L}{\hat{\lambda}_R}+\sigma,
  \]
  连续精确特征值为
  \[
    \lambda_{\mathrm{exact}}=\xi^2+\eta^2+\sigma.
  \]
  由
  \[
    \cos(\xi h)
    =
    1-\frac{\xi^2h^2}{2}
    +\frac{\xi^4h^4}{24}
    -\frac{\xi^6h^6}{720}
    +O(h^8),
  \]
  以及 $\alpha=\xi h,\ \beta=\eta h$，可得
  \begin{equation}
    \hat{\lambda}_R
    =
    \frac{2}{3}
    +
    \frac{1}{6}(\cos\alpha+\cos\beta)
    =
    1-\frac{\xi^2+\eta^2}{12}h^2
    +
    \frac{\xi^4+\eta^4}{144}h^4
    +
    O(h^6).
    \label{eq:lam_R_expand}
  \end{equation}

  另一方面，
  \[
    \hat{\lambda}_L
    =
    \frac{-20+8\cos\alpha+8\cos\beta+4\cos\alpha\cos\beta}{6h^2}.
  \]
  代入 Taylor 展开并整理，得到
  \begin{equation}
    \hat{\lambda}_L
    =
    -(\xi^2+\eta^2)
    +
    \frac{(\xi^2+\eta^2)^2}{12}h^2
    -
    \frac{(\xi^2+\eta^2)(\xi^4+4\xi^2\eta^2+\eta^4)}{360}h^4
    +
    O(h^6).
    \label{eq:lam_L_expand}
  \end{equation}
  由于 $L_h\approx\Delta$，上式的 leading term 为 $-(\xi^2+\eta^2)$，
  与连续 Laplacian 的 Fourier symbol 一致。

  记
  \[
    A=\xi^2+\eta^2,\qquad B=\xi^4+\eta^4.
  \]
  由
  \[
    \hat{\lambda}_R
    =
    1-\frac{A}{12}h^2+\frac{B}{144}h^4+O(h^6),
  \]
  可得
  \[
    \frac{1}{\hat{\lambda}_R}
    =
    1+\frac{A}{12}h^2+\frac{A^2-B}{144}h^4+O(h^6).
  \]
  同时
  \[
    -\hat{\lambda}_L
    =
    A-\frac{A^2}{12}h^2
    +
    \frac{A(\xi^4+4\xi^2\eta^2+\eta^4)}{360}h^4
    +
    O(h^6).
  \]
  因此
  \begin{equation}
    -\frac{\hat{\lambda}_L}{\hat{\lambda}_R}
    =
    \left[
    A-\frac{A^2}{12}h^2+O(h^4)
    \right]
    \left[
    1+\frac{A}{12}h^2+O(h^4)
    \right].
  \end{equation}
  其中 $h^2$ 项为
  \[
    A\cdot\frac{A}{12}-\frac{A^2}{12}=0,
  \]
  故二阶误差项完全消去。继续保留 $h^4$ 项并整理，得到
  \begin{equation}
    -\frac{\hat{\lambda}_L}{\hat{\lambda}_R}
    =
    \xi^2+\eta^2
    -
    \frac{(\xi^2+\eta^2)(3\xi^4-8\xi^2\eta^2+3\eta^4)}{720}h^4
    +
    O(h^6).
  \end{equation}
  加上 $\sigma$ 后即得
  \begin{equation}
    \lambda_{\mathrm{discrete}}(\alpha,\beta)
    =
    \xi^2+\eta^2+\sigma
    -
    \frac{(\xi^2+\eta^2)(3\xi^4-8\xi^2\eta^2+3\eta^4)}{720}h^4
    +
    O(h^6).
  \end{equation}
  这证明了 FFT9 格式在内点 Fourier symbol 意义下具有 $O(h^4)$ consistency。$\qed$
\end{theorem}

\begin{remark}
  定理~\ref{thm:4th_order} 是局部 Fourier symbol 层面的内点一致性分析。
  它说明在周期/局部 Fourier 模式意义下，FFT9 的有效离散特征值
  对连续算子 $-\Delta+\sigma$ 的特征值具有四阶逼近精度。
  该结论本身并不是完整 Dirichlet 边界问题的全局误差估计。
  完整的全局四阶误差还需要边界修正项与九点紧致模板一致，
  并需要离散算子满足相应的稳定性或可逆性条件；
  对于 true Helmholtz，还需避开离散共振或近奇异参数。
\end{remark}

\begin{remark}
  在 Fourier symbol 层面，对比五点差分格式的 leading error
  $\tau^{(5)} = -\frac{\xi^4+\eta^4}{12}h^2 + O(h^4)$，
  九点紧致格式将内点一致性从 $O(h^2)$ 提升到 $O(h^4)$，
  且计算模板仅从5点扩展到9点。
  四阶一致性提升的物理机制在于：$R_h$ 算子的 $O(h^2)$ 误差（含全双调和算子 $\nabla^4$）
  与 $L_h^{(5)}$ 的 $O(h^2)$ 误差（仅含分离的四阶导数）
  的差恰好为混合偏导数 $\partial^4/\partial x^2\partial y^2$，
  该项可被四个角节点的差分模板精确补偿。
  当前 Dirichlet 边界实现的整体四阶收敛行为由第~\ref{sec:experiments}章的光滑 manufactured solution 实验验证。
\end{remark}

\subsubsection{三参数修正族的六阶 local consistency 不可达性}

下面证明，对格式~\eqref{eq:compact-9pt} 进行一种三参数自然推广时，
若参数是与 Fourier mode 无关的统一常数，则无法在内点 Fourier symbol 意义下达到六阶一致性。
该结论仅针对本节定义的三参数九点修正族，而不是所有九点格式或完整 Dirichlet 边值问题的全局误差定理。

\begin{theorem}[三参数九点修正族的六阶 local consistency 不可达性\cite{Sutmann2007}]
  \label{thm:6th_impossible}
  考虑修正紧致格式：
  \begin{equation}
    (-L_h - \alpha h^2 L_h^{(5)}) u + \sigma(R_h + \beta h^2 I) u = (R_h + \gamma h^2 I)f
    \label{eq:modified_compact}
  \end{equation}
  其中 $\alpha, \beta, \gamma \in \mathbb{R}$ 为与 Fourier mode 无关的统一常数。
  若要求该三参数族在 interior Fourier symbol 意义下达到六阶一致性，则展开误差中的
  $h^2$ 项系数
  \[
    c_2(\xi,\eta;\alpha,\beta,\gamma)
  \]
  和 $h^4$ 项系数
  \[
    c_4(\xi,\eta;\alpha,\beta,\gamma)
  \]
  必须对所有连续波数 $(\xi,\eta)$ 同时为零。然而不存在这样的统一常数
  $(\alpha,\beta,\gamma)$。

  \textbf{证明}：对修正格式做 local Fourier symbol 分析，有效离散特征值为：
  \begin{equation}
    \lambda_{\mathrm{mod}} = \frac{-\hat{\lambda}_L - \alpha h^2\hat{\lambda}_5
    + \sigma(\hat{\lambda}_R + \beta h^2)}{\hat{\lambda}_R + \gamma h^2}.
  \end{equation}
  其中 $\hat{\lambda}_5$ 为五点正 Laplacian（近似 $-\Delta$）的 Fourier symbol。
  将 $\lambda_{\mathrm{mod}}-\lambda_{\mathrm{exact}}$ 展开为 $h$ 的幂级数，其 $h^2$ 项系数为
  \begin{equation}
    c_2 = -\alpha(\xi^2+\eta^2) + \beta\sigma - \gamma(\xi^2+\eta^2+\sigma)
    \label{eq:c2_general}
  \end{equation}
  由于六阶 local consistency 至少要求 $c_2=0$ 对所有 $(\xi,\eta)$ 成立，
  当 $\sigma\neq 0$ 时有
  \begin{equation}
    c_2 = (-\alpha-\gamma)(\xi^2+\eta^2) + (\beta-\gamma)\sigma = 0, \quad \forall\,(\xi,\eta),
  \end{equation}
  从而得到必要条件
  \begin{equation}
    \alpha=-\gamma,\qquad \beta=\gamma.
    \label{eq:c2_helmholtz_conditions}
  \end{equation}
  当 $\sigma=0$ 时，$c_2=0$ 同样给出必要条件 $\alpha=-\gamma$。

  将上述必要条件代入 $h^4$ 项后，$c_4$ 仍包含无法由单一常数 $\gamma$ 对所有波数同时消去的
  mode-dependent 项。以 Poisson 情形为例，代入 $\alpha=-\gamma$ 后，$c_4$ 的分子可写为
  \begin{equation}
    N(\xi,\eta,\gamma)
    =
    3\eta^6
    +60\gamma\eta^4
    -5\eta^4\xi^2
    -5\eta^2\xi^4
    +60\gamma\xi^4
    +3\xi^6.
  \end{equation}
  若要求六阶 local consistency，则必须有 $N(\xi,\eta,\gamma)=0$ 对所有 $(\xi,\eta)$ 成立。
  取两个测试波数：
  \[
    N(1,0,\gamma)=60\gamma+3,\qquad
    N(1,1,\gamma)=120\gamma-4.
  \]
  于是 $N(1,0,\gamma)=0$ 给出 $\gamma=-1/20$，
  而 $N(1,1,\gamma)=0$ 给出 $\gamma=1/30$，二者互相矛盾。
  等价地，
  \[
    2N(1,0,\gamma)-N(1,1,\gamma)
    =
    2(60\gamma+3)-(120\gamma-4)
    =
    10\neq 0.
  \]
  因此不存在统一常数 $\gamma$ 使 $c_4$ 对所有波数同时为零。

  对 $\sigma\neq 0$ 的 Helmholtz 情形，先由 $c_2=0$ 得到
  $\alpha=-\gamma,\ \beta=\gamma$，再将其代入 $c_4$；同样取
  $(\xi,\eta)=(1,0)$ 与 $(1,1)$ 会得到互不相容的 $\gamma$ 条件。
  因此，该三参数九点修正族在 Poisson 与 Helmholtz 情形下均无法通过
  mode-independent constants 达到六阶 local Fourier symbol consistency。

  这里的 $(\xi,\eta)$ 是 local Fourier symbol 分析中的连续波数测试，
  不是 Dirichlet DST-I 的离散模态编号，因此在反证中允许取
  $(\xi,\eta)=(1,0)$ 这样的局部波数。该选择只用于展示多项式恒等式不可能对所有波数成立；
  若一个多项式恒等式要求对所有 $(\xi,\eta)$ 成立，则任意两个合法的连续测试点给出矛盾即可否定它。

  Lean 4 辅助验证覆盖的是上述反证中的代数核心：即 $c_2=0$ 推出的参数必要关系，
  以及代入该关系后两个 Fourier 测试模式的 $c_4=0$ 条件互相矛盾。
  具体定理包括 \texttt{poisson\_no\_sixth\_order\_real} 与
  \texttt{helmholtz\_no\_sixth\_order\_real}，验证范围见附录~\ref{app:lean4}。$\qed$
\end{theorem}

\begin{remark}
  定理~\ref{thm:6th_impossible}的结论针对的是三参数修正族
  $(-L_h - \alpha h^2 L_h^{(5)})u + \sigma(R_h+\beta h^2 I)u = (R_h+\gamma h^2 I)f$，
  其中 $\alpha,\beta,\gamma$ 是与 Fourier mode 无关的统一常数。
  该结论是在 local Fourier symbol / interior consistency 意义下成立的，
  不是完整 Dirichlet 边值问题的全局误差不可达结论。
  因而更准确的表述是：在该三参数九点修正族、mode-independent constants
  和 local Fourier symbol 意义下，四阶可视为该族的自然精度上限。
  若允许使用更多参数或更大的计算模板（如引入高阶交叉导数项的差分近似），
  则六阶精度在 Poisson 方程情形下是可达的\cite{Sutmann2007}。
  但这类方案不属于本定理讨论的三参数九点修正族。
  Sutmann\cite{Sutmann2007}通过引入宽度为 $\sqrt{6}h$ 的广义节点，
  构造了六阶精度的紧致差分格式，但该格式的频域除数不再具有简单的解析形式，
  且在 $\sigma \neq 0$ 时仍需额外修正。
\end{remark}

\subsection{Dirichlet 边界条件处理}

对于齐次 Dirichlet 边界条件 $u|_{\partial\Omega} = 0$，
未知数仅在内节点 $(i,j)$，$i,j = 1, \ldots, n-2$ 上取值；
由于边界值为零，边界贡献不改变右端项。

对于非齐次 Dirichlet 边界条件 $u|_{\partial\Omega}=g_D$，
边界值直接赋给边界节点，内节点方程中的已知边界贡献移至右端项。
以靠近左边界的内点 $(1,j)$ 为例，若 $u_{0,j}=g_{0,j}$，则五点格式为
\[
\frac{4u_{1,j}-g_{0,j}-u_{2,j}-u_{1,j-1}-u_{1,j+1}}{h^2}
+\sigma u_{1,j}
=
f_{1,j}.
\]
将未知内点项保留在左端，并将已知边界值移至右端，得到
\[
\frac{4u_{1,j}-u_{2,j}-u_{1,j-1}-u_{1,j+1}}{h^2}
+\sigma u_{1,j}
=
f_{1,j}
+
\frac{g_{0,j}}{h^2}.
\]
因此，非齐次 Dirichlet 边界值在五点正 Laplacian 离散右端项中的修正符号为正号。一般地，
\begin{equation}
  \tilde{f}_{i,j} = f_{i,j} + \frac{1}{h^2}(g_{0,j}\delta_{i,1} + g_{n-1,j}\delta_{i,n-2}
  + g_{i,0}\delta_{j,1} + g_{i,n-1}\delta_{j,n-2})
  \label{eq:nonhom_dirichlet_rhs}
\end{equation}
其中每个与内点相邻的已知边界值都以 $+g/h^2$ 加入右端项。

\subsection{Neumann 边界条件处理}

\subsubsection{Ghost-Point 方法}

对于 Neumann 边界条件 $\partial u/\partial n = 0$，
采用 ghost-point 方法：在边界外引入虚拟节点，
利用中心差分近似法向导数条件：
\begin{equation}
  \frac{u_{-1,j} - u_{1,j}}{2h} = 0 \quad \Longrightarrow \quad u_{-1,j} = u_{1,j}
\end{equation}

将 ghost-point 值代入边界处 $i=0$ 的五点差分格式：
\begin{align}
  \frac{1}{h^2}(4u_{0,j} - u_{-1,j} - u_{1,j} - u_{0,j-1} - u_{0,j+1}) + \sigma u_{0,j} &= f_{0,j} \\
  \Longrightarrow \quad \frac{1}{h^2}(4u_{0,j} - 2u_{1,j} - u_{0,j-1} - u_{0,j+1}) + \sigma u_{0,j} &= f_{0,j}
\end{align}

类似地，对四个边界分别处理。注意 ghost-point 的引入使得 $n$ 个物理节点
（包含边界）全部成为未知量，而非仅内节点。

\subsubsection{对称化与 DCT-I 变换}

ghost-point Neumann Laplacian 矩阵 $G$ 是\textbf{非对称}的。
具体地，以三个物理节点为例，一维 Neumann Laplacian 矩阵为：
\begin{equation}
  G = \frac{1}{h^2}
  \begin{pmatrix}
    2 & -2 & 0 \\
    -1 & 2 & -1 \\
    0 & -2 & 2
  \end{pmatrix}
\end{equation}
可见 $G_{0,1} = -2/h^2 \neq G_{1,0} = -1/h^2$。
同时 $G\mathbf{1}=0$，说明常数向量对应纯 Neumann Poisson 问题的零模态。

引入对角缩放矩阵 $D = \mathrm{diag}(\sqrt{2}, 1, \ldots, 1, \sqrt{2})$，
定义对称化矩阵 $S = D^{-1}GD$。
对三个物理节点，上式给出
\begin{equation}
  S = \frac{1}{h^2}
  \begin{pmatrix}
    2 & -\sqrt{2} & 0 \\
    -\sqrt{2} & 2 & -\sqrt{2} \\
    0 & -\sqrt{2} & 2
  \end{pmatrix},
\end{equation}
因此 $S$ 为对称矩阵。若原始一维系统写作 $Gu=f$，则缩放变量应定义为
\begin{equation}
  \tilde{u}=D^{-1}u,\qquad \tilde{f}=D^{-1}f,
  \qquad S\tilde{u}=\tilde{f},
\end{equation}
求得 $\tilde{u}$ 后再通过 $u=D\tilde{u}$ 恢复原变量。

\begin{theorem}[对称化 DCT-I 对角化\cite{Strang1999}]
  \label{thm:symmetrize}
  对称化矩阵 $S$ 可被 DCT-I 正交对角化。

  \textbf{证明}：对 $S = D^{-1}GD$，考虑其元素。
  以 $S_{0,1} = D^{-1}_{0,0} G_{0,1} D_{1,1} = \frac{1}{\sqrt{2}} \cdot (-\frac{2}{h^2}) \cdot 1 = -\frac{\sqrt{2}}{h^2}$，
  而 $S_{1,0} = D^{-1}_{1,1} G_{1,0} D_{0,0} = 1 \cdot (-\frac{1}{h^2}) \cdot \sqrt{2} = -\frac{\sqrt{2}}{h^2}$。
  因此 $S_{0,1} = S_{1,0}$，$S$ 确实对称。

  定义 DCT-I 正交矩阵 $C \in \mathbb{R}^{n \times n}$，其元素为：
  \begin{equation}
    C_{k,j} = \alpha_k \alpha_j \sqrt{\frac{2}{n-1}} \cos\frac{kj\pi}{n-1}, \quad k,j = 0, \ldots, n-1
    \label{eq:dct1_def}
  \end{equation}
  其中 $\alpha_0 = \alpha_{n-1} = 1/\sqrt{2}$，$\alpha_j = 1$（$1 \leq j \leq n-2$）。
  双端点权重 $\alpha_k \alpha_j$ 保证了 $C$ 的正交性（$C^T C = I$）。

  直接验证 $S = C^T \Lambda C$，其中：
  \begin{equation}
    \lambda_k = \frac{2}{h^2}(1 - \cos\frac{k\pi}{n-1}), \quad k = 0, 1, \ldots, n-1
  \end{equation}
  $k=0$ 对应 $\lambda_0 = 0$（常数模式），这是 Neumann 问题的零特征值。$\qed$

  二维情形中，若 $U$ 的行、列分别对应 $x,y$ 方向，
  标准矩阵写法中右乘方向需使用 $y$ 方向矩阵的转置：
  \[
    G_xU+UG_y^T+\sigma U=F.
  \]
  令
  \[
    \tilde{U}=D_x^{-1}UD_y^{-1},\qquad
    \tilde{F}=D_x^{-1}FD_y^{-1},
  \]
  并记
  \[
    S_x=D_x^{-1}G_xD_x,\qquad
    S_y=D_y^{-1}G_yD_y.
  \]
  由于
  \[
    D_yG_y^TD_y^{-1}
    =
    (D_y^{-1}G_yD_y)^T
    =
    S_y^T
    =
    S_y,
  \]
  得到对称化系统
  \[
    S_x\tilde{U}+\tilde{U}S_y+\sigma\tilde{U}=\tilde{F}.
  \]
  因此 DCT-I 求解流程为：
  \begin{enumerate}[itemsep=1pt]
    \item 缩放右端：$\tilde{F}=D_x^{-1}FD_y^{-1}$
    \item 2D DCT-I 变换：$\widehat{\tilde{F}} = C_x \tilde{F} C_y^T$
    \item 频域求解：
    \[
      \widehat{\tilde{U}}_{k,l}
      = \frac{\widehat{\tilde{F}}_{k,l}}{\lambda_k + \lambda_l + \sigma}
    \]
    \item 逆变换：$\tilde{U} = C_x^T \widehat{\tilde{U}} C_y$
    \item 反缩放：$U = D_x \tilde{U} D_y$
  \end{enumerate}
\end{theorem}
